\documentclass[11pt]{article}

\usepackage[a4paper,margin=1in]{geometry}
\usepackage{amsmath,amssymb,bm}
\usepackage{graphicx}
\usepackage{booktabs}
\usepackage{multirow}
\usepackage{xcolor}
\usepackage{hyperref}
\usepackage[nameinlink,capitalise]{cleveref}
\usepackage[numbers,sort&compress]{natbib}
\usepackage{lineno}
\usepackage{microtype}

% Paper-facing names for the input homotopy.  Implementation identifiers are
% intentionally kept out of the scientific narrative.
\newcommand{\xoff}{\mathcal{X}_{\mathrm{off}}}
\newcommand{\xalign}{\mathcal{X}_{\mathrm{align}}}
\newcommand{\xhlt}{\mathcal{X}_{\mathrm{HLT}}}
\newcommand{\xsupp}[1]{\mathcal{X}^{S}_{#1}}
\newcommand{\xfield}[1]{\mathcal{X}^{F}_{#1}}

\linenumbers

\title{Bridging Offline and Trigger-Level Jet Representations
       through Homotopy-Guided Knowledge Distillation}

\author{
  Ryan Yang \\
  \small Institution \\
  \small \texttt{email@example.com}
}

\date{\today}

\begin{document}

\maketitle

\begin{abstract}
% Write this last. For now, place 3--5 bullets describing:
% the problem, the proposed method, the evaluation, and the central claim.
\end{abstract}

\section{Introduction}
\label{sec:introduction}

\section{Related Work}
\label{sec:related-work}

\section{Problem Formulation}
\label{sec:problem}

\subsection{Paired offline and trigger-level views}

We consider jet classification with paired offline and trigger-level views.
For jet $i$, the offline input $x_i^{\mathrm{off}}\in\xoff$ is constructed
from the richer particle information available after full event
reconstruction, whereas $x_i^{\mathrm{hlt}}\in\xhlt$ contains only the
particles and measurements available to the high-level trigger (HLT).  The
two particle collections need not have the same cardinality, and no
particle-level correspondence is assumed a priori.  They nevertheless share
a canonical jet identity and class label $y_i$, giving the paired sample
\begin{equation}
    \mathcal{D}
    = \left\{
        \left(i,x_i^{\mathrm{off}},x_i^{\mathrm{hlt}},y_i\right)
      \right\}_{i=1}^{N}.
    \label{eq:paired-data}
\end{equation}
Jet-level identity alignment is sufficient for cross-view distillation: a
prediction obtained from one representation can supervise a model evaluated
on another representation of the same physical jet.

The offline view is information-rich but unavailable at trigger inference.
The HLT view is deployable, but some offline particles are absent and the
measurements of corresponding particles may differ because of online
reconstruction.  We use the same classifier family
$g_{\theta}:\mathcal{X}\rightarrow\mathbb{R}^{C}$ throughout, with
$g_{\theta}(x)$ denoting logits over $C$ jet classes.  Models at different
views therefore differ in their available input information, not in their
prediction task or architectural capacity.

A model trained on $\xoff$ defines the projected-offline reference.  An
independently trained, cross-entropy-only model on $\xhlt$ defines the
deployable HLT baseline.  Our objective is to transfer as much of the
reference model's predictive structure as possible to a terminal model whose
inference graph consumes exactly $\xhlt$ and no offline quantities.

\subsection{The limitation of a direct offline-to-HLT transfer}

The most direct knowledge-distillation (KD) strategy trains an exact-HLT
student from the logits of the projected-offline teacher:
\begin{equation}
    g_{\theta_{\mathrm{off}}}(x^{\mathrm{off}})
    \longrightarrow
    g_{\theta_{\mathrm{hlt}}}(x^{\mathrm{hlt}}).
    \label{eq:direct-kd}
\end{equation}
Although the paired jet identities make this objective well-defined, the
student must cross the complete input-information gap in one optimization
problem.  The offline teacher may base its prediction on particles that are
absent online or on particle measurements and identities that differ at HLT.
Its soft distribution can reveal class similarities and relative confidence,
but it cannot make the missing evidence observable to the student.  The
resulting target may therefore be informative yet poorly attainable from the
student's input.

This differs from conventional model compression, in which teacher and
student commonly observe the same input while differing in capacity.  Here,
the architecture remains fixed and deployability is constrained by the input
representation.  We instead seek a continuation in \emph{input information}:
a sequence of paired views connecting the offline and HLT endpoints while
holding the task and classifier family fixed.

Let
\begin{equation}
    x_i^{(0)},x_i^{(1)},\ldots,x_i^{(K)}
    \label{eq:view-sequence}
\end{equation}
be such a sequence, with $x_i^{(0)}=x_i^{\mathrm{off}}$ and
$x_i^{(K)}=x_i^{\mathrm{hlt}}$.  At edge $k$, a teacher evaluated on
$x_i^{(k-1)}$ supervises a freshly initialized student evaluated on
$x_i^{(k)}$.  Our central hypothesis is that nearby teachers define more
attainable targets: each teacher has already adapted to almost the same
observable information as its student while retaining predictive structure
transferred from richer views upstream.  A sequence of such local transfers
may therefore preserve information that is lost in the single projection of
\cref{eq:direct-kd}.  The number and placement of intermediate views determine
the density of this continuation without changing its endpoints.

\section{Homotopy-Guided Knowledge Distillation}
\label{sec:method}

\subsection{A two-coordinate input homotopy}

The offline and HLT representations differ in two conceptually distinct
ways.  First, they may contain different particles.  Second, even when an
offline particle is paired with an HLT particle, their measured features may
differ.  We handle these changes in separate stages rather than modifying
particle support and particle features at the same time.

For jet $i$, define the deterministic view operator
\begin{equation}
    x_i(s,f)
      = V_{s,f}\!\left(x_i^{\mathrm{off}},x_i^{\mathrm{hlt}};A_i\right),
    \qquad s,f\in[0,1],
    \label{eq:view-operator}
\end{equation}
where $A_i$ is a fixed one-to-one particle pairing for one realization of the
homotopy.  The choice of pairing rule is a configurable part of the view
construction rather than a requirement of the distillation objective.  The
support coordinate $s$ controls which particles occupy the input, while the
field coordinate $f$ controls whether paired particles carry offline or HLT
measurements.  The three important views are
\begin{align}
    x_i(0,0) &\in \xoff, \\
    x_i(1,0) &\in \xalign, \\
    x_i(1,1) &= x_i^{\mathrm{hlt}}\in\xhlt.
    \label{eq:view-corners}
\end{align}
The first endpoint, $\xoff$, is the native projected-offline view in the
unified input schema.  Increasing $s$ from zero to one changes its particle
support into the HLT particle support but leaves the features of paired
particles at their offline endpoints.  This produces $\xalign$.  If the HLT
collection is larger than the offline collection, the unavoidable unpaired
HLT particles already carry their native HLT features in $\xalign$.
Increasing $f$ from zero to one then replaces the remaining paired-offline
features with their HLT values, yielding the exact deployable view $\xhlt$.
The path followed in this work is therefore
\begin{equation}
    x_i(s,0):\ \xoff\longrightarrow\xalign,
    \qquad
    x_i(1,f):\ \xalign\longrightarrow\xhlt.
    \label{eq:factorized-homotopy}
\end{equation}

We use \emph{homotopy} in the continuation-method sense: $(s,f)$ indexes an
endpoint-preserving family of related learning problems.  Continuous fields
can vary continuously with these coordinates, but particle insertions,
removals, and categorical changes are necessarily discrete.  For an
individual jet, the resulting tensor path is therefore partly continuous and
partly piecewise constant.  Across the population, the discrete changes are
distributed deterministically and the views remain nested.  We do not assume
that every jet follows a continuous Euclidean trajectory in input space.

\subsection{Full-cardinality bottleneck pairing}

To construct $\xalign$, we must decide which offline and HLT particles refer
to corresponding slots in the homotopy.  Let
$H_i=\{h_{ia}\}_{a=1}^{n_h}$ and
$O_i=\{o_{ib}\}_{b=1}^{n_o}$ be the valid, non-padded HLT and offline
particles in jet $i$.  A one-to-one assignment can contain at most
\begin{equation}
    k_i=\min(n_h,n_o)
    \label{eq:pairing-cardinality}
\end{equation}
pairs.  We always use this maximum cardinality.  Thus, if $n_h\leq n_o$,
every HLT particle receives an offline partner; if $n_h>n_o$, every offline
particle receives an HLT partner.  Only the excess particles in the larger
collection remain unpaired.

Among all maximum-cardinality assignments, we prefer the one that controls
the worst angular separations.  For a candidate pair $(a,b)$, define
\begin{equation}
    \Delta R_{ab}
      = \sqrt{(\eta_{ia}^{\mathrm{hlt}}-\eta_{ib}^{\mathrm{off}})^2
        +(\Delta\phi_{ab})^2},
    \label{eq:pair-delta-r}
\end{equation}
where $\Delta\phi_{ab}$ is the shortest wrapped azimuthal difference.  These
distances are canonically quantized as
$\delta_{ab}=\operatorname{round}_{\mathrm{even}}(\Delta R_{ab}/10^{-7})$
to remove platform-dependent floating-point ambiguity.  For an assignment
$A$, sort its selected distances from largest to smallest:
\begin{equation}
    \bm d(A)
      = \operatorname{sort}_{\downarrow}
        \left(\left\{\delta_{ab}:(a,b)\in A\right\}\right).
    \label{eq:bottleneck-vector}
\end{equation}
We then choose
\begin{equation}
    A_i^{\star}
      = \underset{A:\,|A|=k_i}{\operatorname{arg\,min}_{\mathrm{lex}}}
        \;\bm d(A).
    \label{eq:bottleneck-assignment}
\end{equation}
Lexicographic minimization has a direct interpretation: first minimize the
largest selected $\Delta R$; among assignments attaining that optimum,
minimize the second-largest; then the third-largest; and continue through all
$k_i$ pairs.  Unlike minimizing the mean or sum of $\Delta R$, this objective
does not allow improvements to many easy pairs to compensate for one
avoidable, very poor pair.

Multiple assignments can share the same complete distance vector.  Such ties
are resolved deterministically, in order, using the descending vector of
absolute log transverse-momentum responses, the numbers of particle-category
and valid-charge disagreements, and finally the native offline indices.
These quantities never override a difference in the primary $\Delta R$
objective.

The resulting relation is best understood as a geometric coordinate system
for constructing intermediate views, not as a detector-level truth match.
There is no $\Delta R$, response, category, charge, or confidence threshold,
so full cardinality may require physically imperfect pairs.  The optimization
only guarantees that the ordered tail of selected angular separations is the
best possible among one-to-one assignments of size $k_i$.  We therefore keep
pairing validity distinct from correspondence confidence.  Neither the
assignment indices nor any pairing diagnostic is provided to the classifier.

Full-cardinality bottleneck pairing is the empirically motivated
instantiation studied here, but it is not assumed to be uniquely optimal.
The KD ladder requires only a deterministic, endpoint-consistent pairing from
which nested views can be constructed.  A gated high-purity matcher, a
sum-cost assignment, an optimal-transport construction, or a learned matching
rule could define a different path through the same endpoint representations.
Such alternatives may perform similarly or better, and the appropriate
matching objective should therefore be regarded as a tunable methodological
choice rather than a fixed property of homotopy-guided KD.

Different pairing rules may also induce complementary errors.  One possible
extension is to train a separate ladder for each matching strategy and combine
the predictions of their terminal $\xhlt$ models.  Since every terminal model
receives exact HLT inputs only, this cross-matcher ensemble would not require
offline information at inference.  Where the latency or memory cost of an
ensemble is undesirable, its combined distribution could instead supervise a
single freshly initialized $\xhlt$ model.  We leave both possibilities to
future work; the primary method evaluated here uses one bottleneck pairing
and no model ensemble.

\subsection{Balanced structural transition}

The first leg of the homotopy changes particle support while keeping paired
particle fields at their offline endpoints.  The assignment $A_i^{\star}$
decomposes this change into atomic structural edits.  A paired offline
particle is placed in its corresponding HLT slot; an unused offline particle
is removed; and an unavoidable unpaired HLT particle is inserted with its
native HLT values.  Applying every edit produces exact HLT particle order,
multiplicity, and padding, while paired slots still carry offline fields.  In
other words, the endpoint of the support leg is precisely $\xalign$.

The progress coordinate $s$ determines when these edits occur.  A simple
schedule based only on the number of edits would be poorly calibrated: the
removal of one high-momentum particle can change a jet more than several
low-momentum edits.  We therefore assign every edit $e$ a positive
information-mass proxy $w_e$.  This proxy combines the edit's particle-count
contribution, momentum or energy share, and difference between its endpoint
representations.  A deterministic balancing procedure uses the $w_e$ values
to assign each edit a switch coordinate $\tau_e\in[0,1)$.  The endpoint
convention keeps every edit inactive at $s=0$ and active at $s=1$.

At an interior coordinate $0<s<1$, the support view applies exactly the edits
satisfying
\begin{equation}
    e\ \text{is active at }s
    \quad\Longleftrightarrow\quad
    \tau_e\leq s.
    \label{eq:structural-switch}
\end{equation}
This definition gives three useful properties.  First, the endpoints are
exact: no structural edits are active at $s=0$, and all are active at $s=1$.
Second, the views are nested because each edit switches once and never
reverses.  Third, the construction is reproducible: $\tau_e$ is fixed by the
jet identity, particle pairing, and view configuration, rather than by the
epoch, mini-batch order, or worker process.  Different ladder densities thus
sample the same structural path at different resolutions.

An atomic particle edit cannot be divided continuously.  Individual jets can
therefore still exhibit a finite jump when an influential edit switches.  The
information-mass balancing does not remove this discreteness; it distributes
such changes more evenly across the population and avoids concentrating many
high-impact edits in the same interval of $s$.

\subsection{Uniform matched-field transition}

Once $s=1$, the field path moves every paired particle from its offline field
values to its HLT values.  Unpaired HLT particles are already native and do
not require interpolation.  Define the retained offline strength as
\begin{equation}
    \alpha = 1-f.
    \label{eq:offline-strength}
\end{equation}
For an ordinary continuous field $a$ with offline and HLT endpoints
$a^{\mathrm{off}}$ and $a^{\mathrm{hlt}}$, respectively, the intermediate
value is
\begin{equation}
    a(\alpha)
    = \alpha a^{\mathrm{off}}
      +(1-\alpha)a^{\mathrm{hlt}}.
    \label{eq:continuous-interpolation}
\end{equation}
The same $\alpha$ is applied simultaneously to every paired continuous field.
Angular variables follow the shortest wrapped path, and both endpoints copy
the original offline or HLT values exactly.  The interpolation is not weighted
by a correspondence-confidence score; full-cardinality pairing conveys
validity within the constructed coordinate system, not calibrated match
confidence.

Categorical quantities cannot be linearly interpolated.  Each mismatched
semantic group is instead assigned one deterministic, identity-keyed switch
threshold $r_{i,g}\in[0,1)$.  At offline strength $\alpha$, the group value is
\begin{equation}
    a_{i,g}(\alpha) =
    \begin{cases}
        a_{i,g}^{\mathrm{off}}, & r_{i,g}<\alpha,\\
        a_{i,g}^{\mathrm{hlt}}, & r_{i,g}\geq\alpha.
    \end{cases}
    \label{eq:discrete-switch}
\end{equation}
The threshold is a hash-derived function of the jet identity, HLT particle
slot, semantic group, seed, and view contract; it is never resampled during
training.  Thus, at a view retaining $80\%$ offline strength, approximately
$80\%$ of mismatched groups retain their offline endpoint across the
population, while every individual decision remains reproducible.  Sharing
thresholds across all rungs makes the switches nested: a group changes once
and cannot revert later.  Dependent quantities, including charge,
particle-identity flags, applicability, and validity fields, switch as
coherent atomic groups.  No missing category is invented merely to facilitate
the forced pairing.

We write support-transition views as
$\xsupp{\rho}=V(\rho,0)$, where $\rho$ directly measures progress toward the
target particle support.  Similarly, field-transition views are
$\xfield{\phi}=V(1,\phi)$, where $\phi$ measures progress toward HLT field
values.  The support-aligned pivot satisfies
$\xalign=\xsupp{1}=\xfield{0}$, and the endpoints satisfy
$\xoff=\xsupp{0}$ and $\xhlt=\xfield{1}$.  Both coordinates therefore increase
in the intuitive offline-to-HLT direction.

\subsection{Sequential local distillation}

The main realization uses the following fine-grained path:
\begin{equation}
\begin{split}
    \xoff
      &\rightarrow \xsupp{0.2}\rightarrow \xsupp{0.4}
       \rightarrow \xsupp{0.6}\rightarrow \xsupp{0.8}
       \rightarrow \xalign\\
      &\rightarrow \xfield{0.1}\rightarrow \xfield{0.2}
       \rightarrow \xfield{0.3}\rightarrow \xfield{0.4}
       \rightarrow \xfield{0.5}\\
      &\rightarrow \xfield{0.6}\rightarrow \xfield{0.7}
       \rightarrow \xfield{0.8}\rightarrow \xfield{0.9}
       \rightarrow \xhlt.
    \label{eq:primary-spine}
\end{split}
\end{equation}
The specific grid is not part of the definition of the method.  Coarser or
nonuniform grids select different coordinates from the same $V_{s,f}$ family.

Each arrow in \cref{eq:primary-spine} represents predictive supervision, not
checkpoint continuation.  A student at rung $k$ is initialized from scratch
and receives only the selected teacher distribution from rung $k-1$.  For a
paired jet $i$, let
\begin{align}
    z_i^{(k-1)} &= g_{\theta_{k-1}}(x_i^{(k-1)}),\\
    z_i^{(k)}   &= g_{\theta_k}(x_i^{(k)}).
\end{align}
The softened teacher and student distributions at temperature $T$ are
\begin{equation}
    q_i^{(k-1,T)}=\operatorname{softmax}\!\left(z_i^{(k-1)}/T\right),
    \qquad
    p_i^{(k,T)}=\operatorname{softmax}\!\left(z_i^{(k)}/T\right).
\end{equation}
The student minimizes the per-example objective
\begin{equation}
\begin{split}
    \mathcal{L}_k
      ={}& \lambda_{\mathrm{CE}}
          \operatorname{CE}\!\left(y_i,p_i^{(k,1)}\right) \\
       &+\lambda_{\mathrm{KD}}T^2
          D_{\mathrm{KL}}\!\left(
            q_i^{(k-1,T)}\,\|\,p_i^{(k,T)}
          \right),
    \label{eq:ladder-loss}
\end{split}
\end{equation}
with constant coefficients at every edge.  The class-unweighted hard-label
term prevents the chain from being defined entirely by upstream errors, while
the soft target transfers class similarities and relative confidence learned
from the richer preceding view.

The primary method uses logit distillation only.  Each rung contains one
trained model, and its immediate predecessor is its sole teacher.  No
probability ensemble or representation-level auxiliary objective is inserted
between rungs.  This isolates the effect of traversing intermediate input
views from gains that could instead arise from ensembling or additional
supervision channels.

After training, the checkpoint selected on validation macro-AUC is restored
and evaluated over all training identities to construct the frozen teacher
distribution for the next edge.  Only class probabilities pass down the
spine.  Model parameters, optimizer state, hidden representations, and random
state are not inherited.

This construction replaces one difficult cross-view projection with a
sequence of local projections.  Adjacent teacher and student views share most
of their observable evidence, making the teacher distribution more plausibly
attainable than a target produced directly from $\xoff$.  At the same time,
the teacher carries supervision inherited from richer predecessors.  Greater
path density reduces the nominal input change per edge but also creates more
opportunities to propagate imperfect predictions.  Path density is therefore
an experimental variable, not an assumption that adding rungs must help.

\section{Experimental Setup}
\label{sec:experiments}

\subsection{Data partitions and endpoint construction}

The study uses the complete authenticated population of jets for which the
offline and HLT representations can be paired by canonical identity.  The
registered split contains $2{,}777{,}855$ training jets, $957{,}541$
validation jets, and $899{,}779$ held-out test jets.  Ordinary development,
checkpoint selection, and method comparison access only the training and
validation roles.  The test role remains sealed until the analysis and model
selection procedure are frozen.

% TODO before circulation: describe event generation, detector simulation,
% reconstruction versions, physics-process composition, selections, and the
% exact 15-class label taxonomy.

All homotopy views are derived from the same paired endpoints and immutable
full-cardinality bottleneck assignment.  Assignment integrity requires
exactly $\min(n_h,n_o)$ one-to-one pairs for every nonempty jet, giving
$100\%$ coverage of the smaller particle collection.  Coverage, the selected
$\Delta R$ distribution, its per-jet tail, and comparisons with the gated
high-coverage assignment are retained as matching diagnostics; they are not
classifier inputs or quality cuts.  Neither the class label nor any model
prediction enters view construction.  The homotopy coordinates are likewise
not supplied as model features.  Every classifier must solve the same task
using only the particle information present at its registered view.

\subsection{Registered path-density comparisons}

The fine-grained path in \cref{eq:primary-spine} is the main spine.  Three
additional branches sample the same frozen homotopy more sparsely:
\begin{align}
    \text{direct:}\quad
      &\xoff \rightarrow \xhlt, \\
    \text{coarse:}\quad
      &\xoff \rightarrow \xsupp{0.5} \rightarrow \xalign
       \rightarrow \xfield{1/3} \rightarrow \xfield{2/3}
       \rightarrow \xhlt, \\
    \text{intermediate:}\quad
      &\xoff \rightarrow \xsupp{1/3} \rightarrow \xsupp{2/3}
       \rightarrow \xalign \rightarrow \xfield{0.2}
       \rightarrow \xfield{0.4} \rightarrow \xfield{0.6}
       \rightarrow \xfield{0.8} \rightarrow \xhlt.
    \label{eq:path-density-branches}
\end{align}
Every branch begins from the same authenticated projected-offline teacher.
Models trained at a common coordinate use matched initialization, sampler,
and view seeds across branches, while the causal teacher path changes.  This
design separates the effect of teacher proximity from architecture,
initialization, and endpoint choice.  Every registered branch is evaluated;
validation performance on an early rung does not prune later fits.

\subsection{Model and training protocol}

Every rung uses the same unified Particle Transformer architecture and the
same 21-channel particle-feature schema.  Architectural capacity therefore
cannot explain differences between adjacent rungs, and the terminal $\xhlt$
model has the intended HLT deployment interface.  All models produce logits
over the same 15 classes.  A single authenticated $\xoff$ model is reused as
the common root teacher; every subsequent model is fitted from a fresh
initialization.

Each student is freshly initialized and trained on the full mapped training
population with unweighted cross entropy.  The distillation configuration is
held fixed across the spine:
\begin{equation}
    \lambda_{\mathrm{CE}}=0.25,
    \qquad
    \lambda_{\mathrm{KD}}=0.75,
    \qquad
    T=2.
    \label{eq:registered-loss}
\end{equation}
The mini-batch size is 256.  Optimization uses AdamW with peak learning
rate $3\times10^{-4}$, weight decay $10^{-2}$, $(\beta_1,\beta_2)=(0.9,0.999)$,
and $\epsilon=10^{-8}$.  Forward computation uses bfloat16 where appropriate,
while loss and metric calculations retain the precision required by their
registered implementations.

The learning-rate schedule permits up to 100 complete passes through the
training population:
\begin{itemize}
    \item passes 1--3 linearly warm up to $3\times10^{-4}$;
    \item passes 4--45 hold the peak learning rate;
    \item passes 46--60 use cosine decay to $1.5\times10^{-5}$; and
    \item passes 61--100 retain $1.5\times10^{-5}$ for refinement.
\end{itemize}
Validation is performed after every pass.  Early stopping requires at least
60 passes and uses a patience of 15 passes, with an increase greater than
$5\times10^{-5}$ in macro-AUC required to reset the patience clock.  This
threshold controls stopping only: smaller improvements remain eligible for
exact checkpoint selection.  The selected checkpoint maximizes validation
macro-AUC, with minimum cross entropy, maximum macro log-rejection at $50\%$
signal efficiency, and earliest update used as deterministic tie breakers.

Each fit runs as one process on one GPU with batch size 256.  Teacher
distributions are materialized from the selected checkpoint and joined to
student examples by canonical jet identity.  There is no weight continuation,
multi-model ensemble, representation-level distillation loss, or rolling
optimizer resume.

\subsection{Comparisons and evaluation quantities}

The primary endpoint is the $\xhlt$ model obtained from
\cref{eq:primary-spine}.  Two controls establish the relevant scale.  First,
an exact-HLT model trained with cross entropy alone measures performance
without privileged supervision.  Second, a freshly initialized exact-HLT
student distilled directly from $\xoff$ implements the one-step transfer in
\cref{eq:direct-kd}.  The $\xoff$ model is reported as an information-rich
reference rather than a deployable competitor.  The coarser paths in
\cref{eq:path-density-branches} test whether any advantage of the main method
depends on the proximity of successive teachers.

We report classification accuracy, macro one-vs-rest area under the ROC curve
(AUC), and QCD background rejection at fixed signal efficiency.  To quantify
how much of the offline-to-HLT gap is recovered, let $m_{\mathrm{HLT}}$ be the
metric of the cross-entropy-only exact-HLT baseline, $m_{\mathrm{off}}$ the
metric of $\xoff$, and $m$ the metric of a candidate model.  Its recovery is
defined as
\begin{equation}
    \mathcal{R}_{m}
      = 100\,
        \frac{m-m_{\mathrm{HLT}}}
             {m_{\mathrm{off}}-m_{\mathrm{HLT}}}\,\%.
    \label{eq:metric-recovery}
\end{equation}
For rejection metrics, \cref{eq:metric-recovery} is evaluated in linear
background-rejection space.  Under this convention, the HLT baseline has
$0\%$ recovery and the offline reference has $100\%$ recovery.  Values above
or below that interval are retained rather than clipped.  This normalization
does not replace the underlying metrics; it provides a direct measure of how
much predictive performance lost between the offline and HLT views is restored
by the proposed transfer procedure.

\section{Results}
\label{sec:results}
% Intentionally reserved while the production studies complete.

\section{Discussion}
\label{sec:discussion}

\section{Conclusion}
\label{sec:conclusion}

\bibliographystyle{unsrtnat}
\bibliography{references}

\end{document}
